%==============================================================================
%  CAPÍTULO — EFECTOS LABORALES DE LA APERTURA DEL CONCURSO
%==============================================================================
\chapter{Efectos de la apertura del concurso sobre los contratos de trabajo}
\label{cap:laboral-efectos}

\fichaCapitulo{El concurso como causal de término.}{El artículo 163 bis del Código del
Trabajo, la determinación de la fecha del despido, la suerte de los fueros y la posición
del liquidador como empleador.}

%==============================================================================
\bloqueTeorico
%==============================================================================

\subsection{Dos ordenamientos con lógicas opuestas}

El derecho laboral se construye sobre la estabilidad del vínculo y la protección del
trabajador frente al poder del empleador. El derecho concursal se construye sobre la
liquidación ordenada de un patrimonio insuficiente y la distribución paritaria de la
pérdida. Cuando una empresa con trabajadores entra en liquidación, ambos ordenamientos se
encuentran, y el punto de encuentro es el \artcdt{163 bis} \citeyearpar{codigo-trabajo}.

La solución del legislador es tajante: la resolución de liquidación pone término a los
contratos de trabajo por el solo ministerio de la ley. No hay despido en sentido
estricto, sino terminación legal. Esa calificación tiene consecuencias que recorren todo
este capítulo, y que conviene anticipar:

\begin{enumerate}
  \item No se requiere invocar una causal de las del \artcdt{159}, \artcdt{160} o
  \artcdt{161}: la causal es la resolución misma.
  \item La fecha del término es la de dictación de la resolución de liquidación, no la de
  su notificación ni la de la comunicación al trabajador.
  \item Quien debe cumplir los deberes del empleador ---comunicación, finiquito, pago---
  es el liquidador, con cargo a la masa.
\end{enumerate}

\subsection{Por qué el legislador optó por la terminación automática}

La alternativa ---mantener vigentes los contratos hasta que el liquidador decidiera
individualmente sobre cada uno--- habría producido tres efectos indeseables: la
acumulación de remuneraciones devengadas durante la liquidación como créditos de la masa,
la incertidumbre prolongada del trabajador respecto de su situación, y la imposibilidad
práctica de realizar el activo mientras subsistieran obligaciones laborales corrientes.

La terminación automática corta ese nudo, pero traslada el problema al terreno del pago:
los créditos laborales devengados hasta la fecha del término deben pagarse con
preferencia, y esa preferencia es el objeto del capítulo \ref{cap:creditos-laborales}.

%==============================================================================
\bloqueNormativo
%==============================================================================

\subsection{La causal del artículo 163 bis}

\subsubsection{Supuesto de hecho}

La causal opera con la dictación de la resolución de liquidación respecto del empleador.
No se extiende a la resolución de reorganización: durante la reorganización los contratos
de trabajo subsisten, y su terminación requiere causal propia.

\begin{alerta}[Distinción esencial]
Reorganización y liquidación producen efectos laborales opuestos. La reorganización
mantiene los contratos vigentes; la liquidación los termina de pleno derecho. Confundirlas
en la asesoría genera despidos injustificados o, a la inversa, expectativas de
continuidad que no existen.
\end{alerta}

\subsubsection{Fecha del término}

La ley fija como fecha del término la de dictación de la resolución. La precisión importa
porque de ella dependen el cómputo de la antigüedad, la determinación de las
remuneraciones devengadas y el corte entre créditos concursales y créditos de la
masa.\verificar{Confirmar la redacción vigente del artículo 163 bis en cuanto a la fecha
del término y el cómputo de la indemnización.}

\subsubsection{Indemnizaciones que procede pagar}

\begin{description}[leftmargin=0pt,style=nextline,font=\sffamily\bfseries]
  \item[Indemnización sustitutiva del aviso previo] Procede por mandato legal, con el tope
  que la norma establece.\verificar{Confirmar el tope vigente expresado en unidades de
  fomento.}

  \item[Indemnización por años de servicio] Con los topes generales del Código del
  Trabajo en cuanto a años computables y base de cálculo.

  \item[Feriado proporcional y feriado pendiente] Se devengan hasta la fecha del término.

  \item[Remuneraciones y cotizaciones adeudadas] Con la preferencia que corresponda según
  el capítulo \ref{cap:creditos-laborales}.
\end{description}

\subsection{Los fueros laborales frente a la liquidación}

Es la materia de mayor conflictividad práctica del capítulo. La pregunta es si el fuero
maternal, el fuero sindical y los demás fueros especiales resisten la terminación del
\artcdt{163 bis}, o si ceden ante ella.

La posición mayoritaria sostiene que la causal opera por el solo ministerio de la ley y
no admite excepción por fuero, precisamente porque no se trata de un despido decidido por
el empleador y, por tanto, no hay voluntad que el fuero pueda neutralizar. La posición
contraria invoca el carácter tutelar del fuero y su fundamento constitucional.
\verificar{Sistematizar la jurisprudencia de unificación de la Corte Suprema sobre
subsistencia del fuero maternal y sindical frente al artículo 163 bis, e identificar si
existe criterio asentado.}

\begin{foco}[Cómo tratar el fuero en la práctica]
Mientras la cuestión no esté zanjada, la conducta prudente del liquidador consiste en
documentar la terminación con especial cuidado respecto de los trabajadores aforados,
comunicar a la Inspección del Trabajo y, si el trabajador reclama, provisionar el
resultado del juicio antes de repartir. Repartir el activo sin provisionar una
contingencia de fuero es una fuente clásica de responsabilidad del liquidador.
\end{foco}

\subsection{El liquidador en posición de empleador}

El liquidador no sucede al deudor en la calidad de empleador para todos los efectos, pero
asume los deberes de comunicación, emisión de finiquitos y pago. La \superir{} ha
precisado estos deberes mediante instrucciones específicas, sistematizadas en el capítulo
\ref{cap:normativa-superir}.

\subsection{Continuación de las actividades económicas}

Si la junta acuerda la continuidad del giro, la situación laboral se altera: pueden
celebrarse nuevos contratos, cuyos créditos tienen la naturaleza de créditos de la masa y
no de créditos concursales. La distinción determina el orden de pago y debe quedar
documentada desde el primer día.

\begin{alerta}[Créditos de la masa frente a créditos concursales]
Las obligaciones laborales nacidas antes de la resolución de liquidación son créditos
concursales, sujetos a verificación y prelación. Las nacidas después, con ocasión de la
continuidad de giro, son créditos de la masa y se pagan como gasto de administración. La
confusión entre ambas categorías distorsiona íntegramente el reparto.
\end{alerta}

%==============================================================================
\bloquePractico
%==============================================================================

\subsection{Determinación de la nómina laboral}

\begin{checklist}[Antecedentes que el liquidador debe reunir de inmediato]
  \item Contratos de trabajo vigentes y sus anexos.
  \item Libro de remuneraciones de los últimos veinticuatro meses.
  \item Liquidaciones de sueldo del período impago.
  \item Certificado de cotizaciones previsionales y de salud por trabajador.
  \item Certificado de deuda previsional (Previred o equivalente).
  \item Individualización de trabajadores con fuero vigente y su causa.
  \item Contratos con trabajadores a plazo fijo o por obra, con su fecha de término
  pactada.
  \item Registro de asistencia y de feriados utilizados.
  \item Juicios laborales en curso, con tribunal, rol y estado.
\end{checklist}

\subsection{Cuestionario al empleador deudor}

\begin{cuestionario}[Situación laboral de la empresa]
  \pregunta{¿Cuántos trabajadores tiene, con qué antigüedad y bajo qué modalidad
  contractual?}
  {Define la magnitud del pasivo laboral preferente, que se paga antes que cualquier
  acreedor valista.}

  \pregunta{¿Hay remuneraciones o cotizaciones impagas? ¿Desde cuándo?}
  {Las cotizaciones impagas generan responsabilidades autónomas y agravan la posición de
  los administradores.}

  \pregunta{¿Hay trabajadores con fuero maternal, sindical o de otra clase?}
  {Es la principal contingencia del capítulo y debe provisionarse antes de repartir.}

  \pregunta{¿Existe sindicato o instrumento colectivo vigente?}
  {Los instrumentos colectivos pueden contener indemnizaciones convencionales superiores
  a las legales.}

  \pregunta{¿Hay juicios laborales en curso o demandas anunciadas?}
  {Son créditos contingentes que deben reservarse.}

  \pregunta{¿Se pactaron indemnizaciones convencionales por sobre el mínimo legal?}
  {Su tratamiento preferente tiene topes propios (capítulo \ref{cap:creditos-laborales}).}
\end{cuestionario}

\subsection{Cronograma de la desvinculación}

\begin{flujograma}{Secuencia laboral desde la resolución de liquidación}{cap21-secuencia}
  \node[hito] (res) {RESOLUCIÓN DE LIQUIDACIÓN\\[2pt]
    \normalfont Fecha de término de los contratos (\artcdt{163 bis})};
  \node[nodo, below=of res] (nomina)
    {DETERMINACIÓN DE LA NÓMINA LABORAL\\[2pt]
     \normalfont Antigüedad, remuneración, fueros y deuda previsional};
  \node[nodo, below=of nomina] (comunica)
    {COMUNICACIÓN DEL TÉRMINO AL TRABAJADOR Y A LA INSPECCIÓN DEL TRABAJO};
  \node[nodo, below=of comunica] (calculo)
    {CÁLCULO INDIVIDUAL DE INDEMNIZACIONES Y HABERES ADEUDADOS};
  \node[nodo, below=of calculo] (finiquito)
    {EMISIÓN DEL FINIQUITO ELECTRÓNICO\\[2pt]
     \normalfont Plataforma de la \dt{}};
  \node[favorable, below=of finiquito] (pago)
    {PAGO CON PREFERENCIA DE PRIMERA CLASE\\[2pt]
     \normalfont Sujeto al estado de realización del activo};

  \draw[flecha] (res) -- (nomina);
  \draw[flecha] (nomina) -- (comunica);
  \draw[flecha] (comunica) -- (calculo);
  \draw[flecha] (calculo) -- (finiquito);
  \draw[flecha] (finiquito) -- (pago);
\end{flujograma}

\verificar{Confirmar el plazo legal para comunicar el término a los trabajadores y a la
Inspección del Trabajo en el procedimiento concursal.}

%==============================================================================
\bloqueErrores
%==============================================================================

\errorfrecuente{Aplicar el artículo 163 bis a una reorganización}
{Se despide sin causal legal, con la consiguiente demanda por despido injustificado
contra una empresa que intenta reorganizarse.}
{Recordar que la causal opera sólo con la resolución de liquidación.}

\errorfrecuente{Fijar como fecha del término la de la comunicación al trabajador}
{Se altera el cómputo de antigüedad y de haberes, y se genera diferencia reclamable.}
{Usar siempre la fecha de dictación de la resolución de liquidación.}

\errorfrecuente{Repartir el activo sin provisionar las contingencias de fuero}
{Si el trabajador aforado obtiene sentencia favorable, no hay fondos y la responsabilidad
alcanza al liquidador.}
{Reservar el monto de las contingencias antes de cualquier reparto.}

\errorfrecuente{Confundir créditos de la masa con créditos concursales en la continuidad de giro}
{Se altera el orden de pago y se perjudica a los acreedores preferentes.}
{Llevar registro separado desde el primer día de la continuidad.}

\errorfrecuente{Emitir finiquitos en papel}
{Se incumplen las instrucciones vigentes de la \superir{} sobre finiquitos electrónicos.}
{Tramitarlos por la plataforma de la \dt{}, conforme al capítulo
\ref{cap:normativa-superir}.}
